\documentclass[aspectratio=169]{beamer}

\usetheme{Madrid}
\usecolortheme{seahorse}

\title{Algorithms \& Complexity \\ (Finally Made Simple)}
\author{CSCE 222 Rewrite}
\date{}

\begin{document}

%------------------------------------------------
\begin{frame}
\titlepage
\end{frame}

%------------------------------------------------
\begin{frame}{First: Breathe. This Is Not Magic.}
You are not bad at this. The slides were.

\bigskip

Today we will:
\begin{itemize}
\item Strip everything to patterns.
\item Define every symbol the first time it appears.
\item Connect algorithms $\rightarrow$ counting $\rightarrow$ growth.
\item Show why it is actually simple.
\end{itemize}

\bigskip

Core idea of the entire lecture:

\[
\textbf{How fast does work grow as input size } n \textbf{ grows?}
\]
\end{frame}

%------------------------------------------------
\begin{frame}{Cheat Sheet (Part 1): Core Definitions}
\begin{columns}[t]
\column{0.5\textwidth}

\textbf{Algorithm}

Finite sequence of well-defined steps that:
\begin{itemize}
\item Takes \textbf{input}
\item Produces \textbf{output}
\end{itemize}

Properties:
\begin{itemize}
\item Input
\item Output
\item Definiteness
\item Correctness
\item Finiteness
\item Effectiveness
\item Generality
\end{itemize}

\column{0.5\textwidth}

\textbf{Time Complexity}

Number of operations executed.

\textbf{Space Complexity}

Memory used.

\textbf{Input Size:} $n$

$n$ = number of elements in list.

\textbf{Asymptotic Analysis}

Study behavior as $n \to \infty$.

Ignore:
\begin{itemize}
\item Constants
\item Lower-order terms
\end{itemize}

\end{columns}
\end{frame}

%------------------------------------------------
\begin{frame}{Cheat Sheet (Part 2): Big-O, Big-$\Omega$, Big-$\Theta$}
\begin{columns}[t]
\column{0.5\textwidth}

\textbf{Big-O}

$f(n)$ is $O(g(n))$ if

\[
\exists C>0,\, k>0 \text{ such that }
\]

\[
f(n) \le C g(n) \text{ for all } n>k
\]

Upper bound.

\bigskip

\textbf{Big-$\Omega$}

\[
f(n) \ge C g(n)
\]

Lower bound.

\column{0.5\textwidth}

\textbf{Big-$\Theta$}

$f(n)$ is $\Theta(g(n))$ if:

\[
f(n) \text{ is } O(g(n))
\]
and
\[
f(n) \text{ is } \Omega(g(n))
\]

Tight bound.

\bigskip

Dominant term rule:

Polynomial:
\[
a_n n^n + \dots \Rightarrow O(n^n)
\]

Highest degree dominates.
\end{columns}
\end{frame}

%------------------------------------------------
\begin{frame}{Cheat Sheet (Part 3): Common Growth Rates}
\begin{columns}[t]
\column{0.5\textwidth}

$O(1)$ Constant

$O(\log n)$ Logarithmic

$O(n)$ Linear

$O(n^2)$ Quadratic

\column{0.5\textwidth}

$O(n^b)$ Polynomial

$O(b^n)$ Exponential

\bigskip

Search:
\begin{itemize}
\item Linear: $O(n)$
\item Binary: $O(\log n)$
\end{itemize}

Sorting:
\begin{itemize}
\item Bubble: $O(n^2)$
\item Insertion: $O(n^2)$
\end{itemize}

\end{columns}
\end{frame}

%------------------------------------------------
\begin{frame}{What Is Really Happening?}
Everything in this lecture reduces to:

\[
\textbf{Count operations as a function of } n.
\]

Where:
\begin{itemize}
\item $n$ = number of elements
\item Operation = assignment, comparison, addition
\end{itemize}

We assume:
\[
\text{Each operation takes 1 unit time}
\]

Why?  

Because constants do not matter asymptotically.
\end{frame}

%------------------------------------------------
\begin{frame}{Maximum Finding: Why It Is $O(n)$}
Input:
\[
a_0, a_1, \dots, a_{n-1}
\]

Where:
\begin{itemize}
\item $a_i$ = element at index $i$
\item $n$ = total elements
\end{itemize}

Loop runs from $i=1$ to $n-1$.

Number of comparisons:
\[
n-1
\]

Therefore:
\[
T(n) = c_1 n + c_2
\]

So:
\[
O(n)
\]
\end{frame}

%------------------------------------------------
\begin{frame}{Linear Search: Pattern}
While loop:
\[
i < n \land x \ne a_i
\]

Worst case:
\begin{itemize}
\item Check every element
\item $n$ comparisons
\end{itemize}

Thus:
\[
T(n) = an + b
\]

So:
\[
O(n)
\]

Pattern:
\[
\text{One full pass } \Rightarrow O(n)
\]
\end{frame}

%------------------------------------------------
\begin{frame}{Binary Search: Why It Is $O(\log n)$}
Each step halves search space.

If size is $n$:

After 1 step:
\[
n/2
\]

After 2 steps:
\[
n/4
\]

After $k$ steps:
\[
\frac{n}{2^k}
\]

Stop when:
\[
\frac{n}{2^k} = 1
\]

Solve:
\[
2^k = n
\Rightarrow k = \log_2 n
\]

That is the entire reason.
\end{frame}

%------------------------------------------------
\begin{frame}{Bubble and Insertion: Why $O(n^2)$}
Outer loop runs $n$ times.

Inner loop runs about $n$ times.

Total operations:
\[
n \times n = n^2
\]

More precisely:
\[
\frac{n(n-1)}{2}
\]

But asymptotically:
\[
O(n^2)
\]

Pattern:
\[
\text{Nested loops over } n \Rightarrow O(n^2)
\]
\end{frame}

%------------------------------------------------
\begin{frame}{Big-O Definition (Fully Unpacked)}
Definition:

\[
f(n) \le C g(n) \text{ for all } n>k
\]

Variables:

\begin{itemize}
\item $f(n)$ = actual runtime
\item $g(n)$ = comparison function
\item $C$ = constant multiplier
\item $k$ = threshold input size
\end{itemize}

Why $C$?

Allows scaling difference.

Without $C$:
\[
4n^2 \not\le n^2
\]

Why $k$?

Small $n$ can behave irregularly.
\end{frame}

%------------------------------------------------
\begin{frame}{Constructing Witness Constants (Example)}
Show:
\[
f(n) = 4n^2 + 1000n + 904
\]
is $O(n^2)$.

We want:
\[
f(n) \le C n^2
\]

For $n \ge 1$:

\[
1000n \le 1000n^2
\]
\[
904 \le 904n^2
\]

So:
\[
f(n) \le (4+1000+904)n^2
\]

\[
C=1908,\quad k=1
\]

Witnesses constructed.
\end{frame}

%------------------------------------------------
\begin{frame}{Why $n^2$ Is NOT $O(n)$}
Assume:
\[
n^2 \le Cn
\]

Divide by $n$:
\[
n \le C
\]

But $n$ grows without bound.

Eventually:
\[
n > C
\]

Contradiction.

Therefore:
\[
n^2 \notin O(n)
\]

Growth rate matters.
\end{frame}

%------------------------------------------------
\begin{frame}{Big-$\Omega$ and Big-$\Theta$ Intuition}
Big-$\Omega$:

\[
f(n) \ge C g(n)
\]

Lower bound.

Big-$\Theta$:

Both:
\[
C_1 g(n) \le f(n) \le C_2 g(n)
\]

Meaning:

$f(n)$ grows at the same rate as $g(n)$.

Example:

\[
4n^2 + 3n \in \Theta(n^2)
\]
\end{frame}

%------------------------------------------------
\begin{frame}{Homework Survival Strategy}
When you see a function:

Step 1: Identify dominant term.

Step 2: Ignore constants.

Step 3: Ignore lower-order terms.

\bigskip

When you see loops:

\begin{itemize}
\item One loop over $n$ $\Rightarrow O(n)$
\item Nested loops $\Rightarrow O(n^2)$
\item Halving each time $\Rightarrow O(\log n)$
\end{itemize}

That covers 90\% of test questions.
\end{frame}

%------------------------------------------------
\begin{frame}{Extra Example (Deep Why)}
Find Big-O of:

\[
f(n) = 9n^2 + 8n\log n
\]

Compare growth:

\[
n^2 \text{ vs } n\log n
\]

Since:
\[
\frac{n^2}{n\log n} = \frac{n}{\log n} \to \infty
\]

$n^2$ grows faster.

So:
\[
O(n^2)
\]
\end{frame}

%------------------------------------------------
\begin{frame}{Final Perspective}
This entire lecture reduces to:

\[
\textbf{How many times does the work repeat?}
\]

Count repetition pattern.

Match to growth pattern.

Ignore noise.

\bigskip

You are not lost.

You were just given too much noise.

Now you see the structure.
\end{frame}

%------------------------------------------------
\end{document}